\section{Verification methodology}
\label{p:verification}

The verification policy underlying this paper is simple to state and
tight in practice: \emph{the build of the manuscript is gated by a test
suite}.

\subsection{Two backends, one policy}
Algebraic identities are checked by \texttt{sympy}~\cite{Meurer2017}:
matrix equalities, finite Taylor expansions, Clifford and Pauli
algebra, eigenvalue computations. Combinatorial and order-theoretic
claims --- Nielsen--Ninomiya balance, isotropy as SMT, bipartite index,
structural consistency of small causal sets --- are checked by the SMT
solver \texttt{z3}~\cite{deMoura2008}. The two backends together cover
the algebraic skeleton on which the physicist's manipulations rest.

\subsection{Pipeline and numbers}
A single build script runs \texttt{pytest} and only proceeds with
\texttt{latexmk} on a clean exit. The numbers are summarized in
Table~\ref{p:tab:verification}.

\begin{table}[h]
\centering
\begin{tabular}{lr}
\toprule
Validator modules                & $38$ \\
Tests (parametric)               & $290$ \\
Pure \texttt{sympy} modules      & $33$ \\
Pure \texttt{z3} modules         & $4$ \\
Mixed \texttt{sympy}+\texttt{z3} & $1$ \\
End-to-end build time            & $\sim 15\text{ s}$ \\
\texttt{\% verified-by:} comments inserted & $163$ \\
Chapters covered (long form)     & $32/32$ \\
\bottomrule
\end{tabular}
\caption{Verification suite at the time of writing. The long-form
companion of this paper is itself built by the same pipeline; the
present manuscript inherits the same policy.}
\label{p:tab:verification}
\end{table}

\subsection{Granularity}
Each section of the long form (and the equivalent places in this paper)
carries inline references of the form
\begin{center}\small\texttt{\% verified-by: validators/\textless module\textgreater.py::\textless test\textgreater}\end{center}
linking each claim to its test. A representative example: the Clifford
anticommutator check, in essence,
\begin{verbatim}
def clifford_holds() -> bool:
    g = dirac_gamma_matrices()
    eta = minkowski_metric()
    I4, Z4 = sp.eye(4), sp.zeros(4, 4)
    for mu in range(4):
        for nu in range(4):
            diff = sp.simplify(
                g[mu] * g[nu] + g[nu] * g[mu]
                - 2 * eta[mu, nu] * I4
            )
            if diff != Z4:
                return False
    return True
\end{verbatim}
runs sixteen anticommutator checks against the Minkowski metric in
under one second. The Lorentz algebra closure
[Eq.~\eqref{p:eq:lorentz-closure}] is verified analogously across all
$256$~quadruples $(\mu,\nu,\rho,\sigma)$.

\subsection{Errors caught during development}
The policy proved useful, not just decorative. Three classes of error
were caught and corrected by the suite during the writing of the
long-form manuscript: (i)~an overconfident combinatorial predicate that
asserted the wrong restriction on chirality assignments
(SMT correctly reported \texttt{unsat}); (ii)~a sign error in the
spinorial covariance identity [Eq.~\eqref{p:eq:cov-id}], caught by the
closed-form symbolic test; (iii)~an inversion in the docstring of the
Hermitian / anti-Hermitian status of rotation vs.\ boost generators,
caught by \texttt{::test\_rotation\_generators\_hermitian} and
\texttt{::test\_boost\_generators\_antihermitian}. In each case, the
fix was algorithmic and immediately reproducible, not interpretive.

\subsection{Comparison with proof assistants}
A small but growing body of work uses interactive theorem provers
(Lean, Coq, Isabelle) to formalize portions of mathematics and
physics~\cite{Buzzard2022}. The present approach is deliberately
coarser. We do not certify the entire physics in a typed proof
assistant; we certify the algebraic and combinatorial skeleton on which
the standard physicist's manipulations rest, leaving steps such as the
infinite-dimensional limit of the continuum as theorems argued by hand.
The trade-off is favourable to the working scientist: the verification
fits naturally inside an existing \texttt{pytest} + \texttt{latexmk}
workflow, runs in tens of seconds, and produces manuscripts whose
algebraic claims are auditable in under a minute per claim by reading
the test that underwrites them.
